%------------------------------------------------------------------------------%
\documentclass[fleqn,utf8,aspectratio=169,12pt]{beamer}

\usepackage{integracao_e_series}

\title{Revisão de Limites}

%------------------------------------------------------------------------------%
\begin{document}

\addtitle

%------------------------------------------------------------------------------%
\section{Limites}

%------------------------------------------------------------------------------%
\begin{frame}{Limite de uma Função}
  Seja $f(x)$ definida em um intervalo aberto em torno de $a$,

  \pause
  exceto, possivelmente, no próprio $a$

  \pause
  O limite de $f(x)$ quando $x$ tende para $a$ é o número $L$
  \[
    \lim_{x\to a} f(x) = L
  \]
  \pause
  se, para cada \(\epsilon > O\), existir \(\delta > O\), tais que
  \pause
  \[
    0 < \abs{x-a} < \delta \qquad\Rightarrow\qquad \abs{f(x)-L} < \epsilon
  \]
\end{frame}

%------------------------------------------------------------------------------%
\begin{frame}{Limites de algumas funções}
  \( \lim_{x\to a} c = c \)

  \pause
  \medskip
  \( \lim_{x\to a} x = a \)

  \pause
  \medskip
  \( \lim_{x\to a} x^n = a^n \)
  \tabto{35mm}
  $n$ número inteiro positivo

  \pause
  \medskip
  \( \lim_{x\to a} \sqrt[n]{x} = \sqrt[n]{a} \)
  \tabto{35mm}
  $n$ número inteiro positivo, se $n$ for par precisamos impor $a\geq 0$

  \pause
  \medskip
  \( \lim_{x\to a} p(x) = p(a) \)
  \tabto{35mm}
  onde $p(x)$ é um polinômio em $x$
\end{frame}

%------------------------------------------------------------------------------%
\begin{frame}{Limites de algumas funções}
  \( \lim_{x\to \infty} \frac{1}{x} = 0 \)

  \pause
  \medskip
  \( \lim_{x\to 0^+} \frac{1}{x} = \infty \)

  \pause
  \medskip
  \( \lim_{x\to 0^-} \frac{1}{x} = -\infty \)

\end{frame}

%------------------------------------------------------------------------------%
\begin{frame}{Propriedades de Limites}
  Se os limites \(\lim_{x\to a}f(x)\) e  \(\lim_{x\to a}g(x)\) existirem, então:\\[4mm]
  \pause
  \begin{enumerate}[<+->]
    \item
      \(
        \lim_{x\to a} \left(f(x)+g(x)\right)
        = \lim_{x\to a}f(x) + \lim_{x\to a}g(x)
      \)
      \hfill Soma
      \\[5mm]

    \item
      \(
        \lim_{x\to a} \left(f(x)-g(x)\right)
        = \lim_{x\to a}f(x) - \lim_{x\to a}g(x)
      \)
      \hfill Diferença
      \\[5mm]

    \item
      \(
        \lim_{x\to a} \left(c f(x)\right)
        = c\lim_{x\to a} f(x),\quad \forall\; c\in\R
      \)
      \hfill Multiplicação por escalar
      \\[5mm]

    \item
      \(
        \lim_{x\to a} \left(f(x)g(x)\right)
        = \left(\lim_{x\to a} f(x)\right)\left(\lim_{x\to a} g(x)\right)
      \)
      \hfill Produto
      \\[5mm]

    \item
      \(
        \lim_{x\to a} \frac{f(x)}{g(x)}
        = \dfrac{\displaystyle \lim_{x\to a} f(x)}{\displaystyle \lim_{x\to a} g(x)}\quad
      \)
      desde que $\;\lim_{x\to a} f(x) \neq 0$
      \hfill Quociente
  \end{enumerate}
\end{frame}

%------------------------------------------------------------------------------%
\begin{frame}{Teorema do confronto}
  Se \(g(x) < j(x) < h(x)\)

  \pause
  para todo $x$ em um intervalo aberto contendo $a$,
  \pause
  exceto $x = a$,
  \pause
  e
  \[
    \lim_{x\to a} g(x) = \lim_{x\to a} h(x) = L
  \]
  \pause
  então
  \[
   \lim_{x\to a} f(x) = L
  \]
\end{frame}

%------------------------------------------------------------------------------%
\section{Técnicas para Calcular Limites em uma Variável}

%------------------------------------------------------------------------------%
\begin{frame}{Estratégia Geral para Cálculo de Limites}
\begin{enumerate}[<+->]
    \item Substituir diretamente o valor da variável no limite
    \item Identificar o resultado: número finito, $\pm\infty$ ou indeterminação
    \item Em caso de indeterminação:
    \begin{itemize}
        \item aplicar manipulações algébricas
        \item utilizar fatorações
        \item racionalizações
        \item dividir pelo termo de maior grau
    \end{itemize}
\end{enumerate}
\end{frame}

%------------------------------------------------------------------------------%
\begin{frame}{Indeterminações Comuns}
\begin{itemize}[<+->]
    \setlength{\itemsep}{5mm}
    \item $\dfrac{0}{0}$           \tabto{35mm} indeterminação algébrica
    \item $\dfrac{\infty}{\infty}$ \tabto{35mm} formas racionais
    \item $\infty - \infty$        \tabto{35mm} diferença de infinitos
    \item $0 \cdot \infty$
    \item $1^\infty$,\; $0^0$,\; $\infty^0$
\end{itemize}
\end{frame}

%------------------------------------------------------------------------------%
%------------------------------------------------------------------------------%
\begin{question}
  Avalie o limite
  \(
    \lim_{x \to 2} \frac{x^2 - 4}{x-2}
  \)
\end{question}

%------------------------------------------------------------------------------%
\begin{resolution}{Solução}
  Substituindo diretamente
  \pause
  \[
    \lim_{x \to 2} \frac{x^2 - 4}{x-2}
    \pause
    \;=\; \frac{x^2 - 4}{x-2} \evalat{x=2}{}
    \pause
    \;=\; \frac{2^2 - 4}{2-2}
    \pause
    \;=\; \frac{0}{0}
  \]

  \pause
  Indeterminação \(\sfrac{0}{0}\),

  \pause
  não podemos dizer nada sobre o limite
\end{resolution}

%------------------------------------------------------------------------------%
\begin{resolution}{Solução}
  \onslide<+->{Fatorando o numerador}
  \begin{align*}
    \onslide<+->{f(x)}
    \onslide<+->{& = \frac{x^2 - 4}{x-2} \\}
    \onslide<+->{& = \frac{(x-2)(x+2)}{x-2} \\}
    \onslide<+->{& = x+2}
  \end{align*}

  \onslide<+->{Calculando o Limite}
  \[
    \onslide<+->{\lim_{x \to 2} f(x)}
    \onslide<+->{\;=\; \lim_{x \to 2} (x+2)}
    \onslide<+->{\;=\; (x+2) \evalat{x=2}{}}
    \onslide<+->{\;=\; 4}
  \]
\end{resolution}

%------------------------------------------------------------------------------%
%------------------------------------------------------------------------------%
\begin{question}
Calcule o limite
\(
  \lim_{x \to \infty} \frac{3x^2 + 5x}{2x^2 - x}
\)
\end{question}

%------------------------------------------------------------------------------%
\begin{resolution}{Solução}
  Não podemos substituir diretamente,
  \pause
  mas sabemos que
  \pause
  \[
    \lim_{x \to \infty} 3x^2 + 5x = \infty
  \]
  \[
    \pause
    \lim_{x \to \infty} 2x^2 - x = \infty
  \]
  \pause
  portanto temos uma indeterminação \(\sfrac{\infty}{\infty}\)

  \pause
  e não podemos dizer nada sobre o limite
\end{resolution}

%------------------------------------------------------------------------------%
\begin{resolution}{Solução}
  Dividindo o numerador e o denominador pelo termo de maior grau, $x^2$
  \pause
  \[
    f(x)
    \pause
    \;=\; \frac{3x^2 + 5x}{2x^2 - x}
    \pause
    \;=\; \frac{\left(3x^2 + 5x\right)\sfrac{}{x^2}}{\left(2x^2 - x\right)\sfrac{}{x^2}}
    \pause
    \;=\; \frac{3 + \sfrac{5}{x}}{2 - \sfrac{1}{x}}
  \]
  \pause
  Sabemos que
  \(
    \lim_{x\to\infty} \frac{1}{x} = 0
  \),
  \pause
  então
  \[
    \pause
    \lim_{x\to\infty} f(x)
    \pause
    \;=\; \lim_{x\to\infty} \frac{3 + \sfrac{5}{x}}{2 - \sfrac{1}{x}}
    \pause
    \;=\; \frac{3 + 0}{2 - 0}
    \pause
    \;=\; \frac{3}{2}
  \]
\end{resolution}

%------------------------------------------------------------------------------%
%------------------------------------------------------------------------------%
\begin{question}
  Avalie
  \(
    \lim_{x \to \infty} \left(\sqrt{x^2+3x} - x\right)
  \)
\end{question}

%------------------------------------------------------------------------------%
\begin{resolution}{Solução}
  Não podemos substituir diretamente,
  \pause
  mas sabemos que
  \pause
  \[
    \lim_{x \to \infty} \sqrt{x^2+3x} = \infty
  \]
  \[
    \pause
    \lim_{x \to \infty} x = \infty
  \]
  \pause
  portanto temos uma indeterminação \(\infty-\infty\)

  \pause
  e não podemos dizer nada sobre o limite
\end{resolution}

%------------------------------------------------------------------------------%
\begin{resolution}{Solução}
\begin{columns}[t]
\column{0.5\linewidth}
  \begin{align*}
    \onslide<+->{f(x)}
    \onslide<+->{& = \sqrt{x^2+3x} - x \\}
    \onslide<+->{& = \left(\sqrt{x^2+3x} - x\right)\; \frac{\sqrt{x^2+3x} + x}{\sqrt{x^2+3x} + x} \\}
    \onslide<+->{& = \frac{x^2+3x - x^2}{\sqrt{x^2+3x} + x} \\}
    \onslide<+->{& = \frac{3x}{\sqrt{x^2+3x} + x} \\}
    \onslide<+->{& = \frac{3x}{\sqrt{x^2\left(1+\dfrac{3}{x}\right)} + x}}
  \end{align*}
\column{0.5\linewidth}
  \begin{align*}
    \onslide<+->{f(x)}
    \onslide<+->{& = \frac{3x}{x\sqrt{1+\dfrac{3}{x}\;} + x} \\}
    \onslide<+->{& = \frac{3x}{x\left(\sqrt{1+\dfrac{3}{x}\;} + 1\right)} \\}
    \onslide<+->{& = \frac{3}{\sqrt{1+\dfrac{3}{x}\;} + 1}}
  \end{align*}
\end{columns}
\end{resolution}

%------------------------------------------------------------------------------%
\begin{resolution}{Solução}
  \begin{align*}
    \onslide<+->{\lim_{x \to \infty} f(x)}
    \onslide<+->{& = \lim_{x \to \infty} \frac{3}{\sqrt{1+\dfrac{3}{x}\;} + 1} \\}
    \onslide<+->{& = \frac{3}{\sqrt{\displaystyle 1+\lim_{x \to \infty}\left(\frac{3}{x}\right)\;} + 1} \\}
    \onslide<+->{& = \frac{3}{\sqrt{1+0\;} + 1}}
    \onslide<+->{\;=\; \frac{3}{2}}
  \end{align*}
\end{resolution}

%------------------------------------------------------------------------------%
%------------------------------------------------------------------------------%
\begin{question}
  Avalie
  \(
    \lim_{x \to 0^+} x\ln(x)
  \)
\end{question}

%------------------------------------------------------------------------------%
\begin{resolution}{Solução}
  Não podemos substituir diretamente,
  \pause
  mas sabemos que
  \[
    \pause
    \lim_{x \to 0^+} x = 0
  \]
  \[
    \pause
    \lim_{x \to 0^+} \ln x = -\infty
  \]
  \pause
  portanto temos uma indeterminação \(0(-\infty)\)

  \pause
  e não podemos dizer nada sobre o limite
\end{resolution}

%------------------------------------------------------------------------------%
\begin{resolution}{Solução}
  \begin{align*}
    \onslide<+->{\lim_{x \to 0^+} x\ln(x)}
    \onslide<+->{& = \lim_{x \to 0^+} \frac{\ln(x)}{\sfrac{1}{x}}  & & \text{indeterminação } \frac{\infty}{\infty}\\}
    \onslide<+->{& = \lim_{x \to 0^+} \frac{\left(\ln(x)\right)'}{\left(x^{-1}\right)'} & & \text{Regra de L'Hôpital} \\}
    \onslide<+->{& = \lim_{x \to 0^+} \frac{x^{-1}}{-x^{-2}} \\}
    \onslide<+->{& = \lim_{x \to 0^+} \frac{x^2}{-x} \\}
    \onslide<+->{& = \lim_{x \to 0^+} -x}
    \onslide<+->{\;=\; 0}
  \end{align*}
\end{resolution}

%------------------------------------------------------------------------------%
%------------------------------------------------------------------------------%
\begin{question}
  Calcule
  \(
    \lim_{x \to 0} x\sin\left(\frac{1}{x}\right)
  \)

  \bigskip
  \pause
  Dica: use o Teorema do Confronto
\end{question}

%------------------------------------------------------------------------------%
\begin{resolution}{Solução}
  Como
  \[
    -1 \leq \sin(u) \leq 1 \qquad \forall u \in \R
  \]
  \pause
  para $x\neq 0$
  \[
    -1 \leq \sin\left(\frac{1}{x}\right) \leq 1
  \]
  \pause
  portanto
  \[
    -\abs{x} \leq x\sin\left(\frac{1}{x}\right) \leq \abs{x}
  \]
\end{resolution}

%------------------------------------------------------------------------------%
\begin{resolution}{Solução}
  Sabemos que
  \[
    \lim_{x\to 0} \abs{x} \;=\; 0
    \pause
    \qquad\text{e}\qquad
    \lim_{x\to 0} -\abs{x} \;=\; 0
  \]

  \pause
  Portanto, pelo Teorema do Confronto,
  \pause
  \[
    \lim_{x \to 0} x\sin\left(\frac{1}{x}\right) = 0
  \]
\end{resolution}

%------------------------------------------------------------------------------%

%------------------------------------------------------------------------------%
\begin{frame}{Técnicas principais}
\begin{itemize}
  \item Substituição direta
  \item Fatoração e simplificação algébrica
  \item Divisão por maior potência para limites no infinito
  \item Racionalização em diferenças envolvendo raízes
  \item Reescrita de produtos em quocientes
  \item Uso da Regra de L'Hôpital
\end{itemize}
\end{frame}

%------------------------------------------------------------------------------%
\section{Lista Mínima}

%------------------------------------------------------------------------------%
\begin{frame}{Lista Mínima}
Apostila ``\href{https://sites.google.com/view/prof-luis-dafonseca/gradua\%C3\%A7\%C3\%A3o/Integrais-e-Series}{Integração e Séries}''
\begin{enumerate}
  \item Estudar atentamente texto do Capítulo 1
\end{enumerate}
\vfill
Atenção: A prova é baseada no livro, não nas apresentações
\end{frame}

%------------------------------------------------------------------------------%
\end{document}
%------------------------------------------------------------------------------%
